\documentclass[11pt]{amsart}
\usepackage{amsmath,amssymb}
\setlength{\emergencystretch}{3em}
\tolerance=600
\usepackage{tikz-cd}
\usepackage[hidelinks]{hyperref}

\newtheorem{thm}{Theorem}[section]
\newtheorem{lem}[thm]{Lemma}
\newtheorem{cor}[thm]{Corollary}
\newtheorem{prop}[thm]{Proposition}
\theoremstyle{definition}
\newtheorem{defi}[thm]{Definition}
\newtheorem{rem}[thm]{Remark}
\newtheorem{ex}[thm]{Example}
\newtheorem{que}[thm]{Question}

\newcommand{\Same}{\mathrm{Same}}
\newcommand{\Fix}{\mathrm{Fix}}
\newcommand{\uco}{\mathrm{uco}}
\newcommand{\Psh}{\mathrm{Psh}}
\newcommand{\Sh}{\mathrm{Sh}}
\newcommand{\Ran}{\mathrm{Ran}}
\newcommand{\cP}{\mathcal P}
\newcommand{\cE}{\mathcal E}

\begin{document}

\title[Abstract domains and cut points in the sameness lattice]{Abstract domains and cut points in the lattice of two-out-of-three classes: joins of closure operators on a lattice and of cut-point topologies on a control-flow graph}
\author{Ryota Shibusawa}
\thanks{Corresponding author. E-mail: \texttt{r-shibusawa@daiichi-koudai.ac.jp}}
\address{Daiichi Institute of Technology}
\email{r-shibusawa@daiichi-koudai.ac.jp}
\keywords{closure operator, abstract interpretation, abstract domain, reduced product, Grothendieck topology, presheaf topos, control-flow graph, cut point, two-out-of-three, join of localisations}
\subjclass[2020]{18F10 (primary), 06A15, 68Q55, 18B25 (secondary)}

\begin{abstract}
An abstract domain of Cousot and Cousot, a closure operator $\rho$ on a
complete lattice, determines the class $W_\rho$ of pairs
$x\le y$ that $\rho$ does not distinguish, an element of the lattice of
two-out-of-three classes. We show that
$W_\rho\cap W_{\rho'}$ is the class of the reduced product, that
$W_\rho\vee W_{\rho'}$ is the class of pairs with $y\le(\rho\rho')^kx$ for
some $k$, and that it is the class of the least common abstraction iff the
alternation $\rho\rho'\rho\rho'\cdots$ stops at a finite stage at every
point; cancellation is needed when the closures do not preserve meets, as
for reachability closures. On the presheaf
topos of the path category of a directed graph, a set $S$ of vertices gives
the topology whose sheaves are the presheaves determined by their values at
$S$. For two such sets the join of the two classes inverted by
sheafification is the class of the join topology iff the alternations
between $S$ and $T$ along paths avoiding $S\cap T$ are bounded; for a
finite graph, iff no strongly connected component of the graph with
$S\cap T$ removed meets both $S$ and $T$, a linear-time test, and then the
join topology is that of $S\cap T$.
\end{abstract}

\maketitle

\section{Introduction}

In one sentence: for two abstractions of the same concrete lattice, the
join of the two classes of ``pairs the abstraction cannot tell apart'' is
generated by alternating the two abstractions and cancelling, and it is
the class of the least common abstraction exactly when the alternation is
finite at every point; on the path category of a control-flow graph, with
abstractions given by sets of cut points, the finiteness is a condition on
closed walks, tested in linear time.

Let $\mathcal C$ be a category. A \emph{sameness} on $\mathcal C$ is a
class of morphisms containing the identities and closed under composition
and two-out-of-three; samenesses form a complete lattice $\Same(\mathcal
C)$, the meet being the intersection and the join the least sameness
containing both \cite{Same}. Every localisation of $\mathcal C$
contributes the class $W$ of maps it inverts, and the question studied in
\cite{Topjoin,Finitejoin,Stabjoin,Fijoin} is whether $W$ sends the join
of two localisations to the join of the two classes. For localisations
given by closure operators on subobjects (Lawvere--Tierney topologies,
hereditary torsion theories, nuclei) the answer of \cite{Fijoin} is a
finite-iteration criterion: in a Grothendieck topos or a Grothendieck
abelian category, $W_c\vee W_{c'}=W_{\bar c}$ for the least closure
operator $\bar c$ above $c$ and $c'$ iff $(cc')^k=\bar c$ for some $k$.

This paper applies the criterion in two settings suggested by the
semantics of programs. The first is the lattice of abstract domains of
abstract interpretation \cite{CC79}: an abstract domain on a complete
lattice $L$ of concrete properties is an upper closure operator
$\rho$ on $L$, the closure operators form a complete lattice
$\uco(L)$ \cite{Ward}, the reduced product of two domains is their meet
and their least common abstraction is their join \cite{GR97}. The
category is the poset $L$ itself, so that the maps are the inequalities
$x\le y$, and the class $W_\rho$ consists of the inequalities that $\rho$
collapses, $\rho x=\rho y$: the refinements $y$ of $x$ that the analysis
cannot see. General abstract domains are not nuclei, they need not
preserve binary meets (intervals, octagons, reachability under a
transition relation), and the criterion of \cite{Fijoin} for nuclei on a
frame does not apply to them as it stands. Theorem~\ref{thm:lattice}
treats arbitrary closure operators on a complete lattice: the meet is
always preserved, the join is the class of pairs $x\le y$ with $y\le
(\rho\rho')^kx$ for some $k$, it is the class of $\rho\sqcup\rho'$ iff the
alternation reaches $\rho\sqcup\rho'$ at a finite stage at every $x$,
and, unlike the case of nuclei, the two-out-of-three cancellation is
needed to generate the join from the alternating chains (already on the
four-element Boolean lattice). Section~\ref{sec:reach} spells this out for
the reachability closures $S\mapsto R^*(S\cup I)$ of transition relations:
the join of the two classes is bounded alternation of the two relations,
it is the class of the union relation iff the alternation depth is
finite from every set of states, which holds for finite state spaces with
bound the number of states and fails on $\mathbb N$ for the even and the
odd steps.

The second setting is the presheaf topos on the path category of a
directed graph $G$, a control-flow graph: a presheaf is a flow of sets
along the control flow, and a set $S$ of vertices, a set of cut points in
the sense of Bourdoncle \cite{Bourdoncle}, determines the Grothendieck
topology $J_S$ whose sheaves are the presheaves determined by their values
at $S$, $\Sh(J_S)\simeq\Psh(S)$. Theorem~\ref{thm:cut} says that
$W_{J_S}\vee W_{J_T}=W_{J_S\vee J_T}$ iff the number of alternations
between visits to $T$ and visits to $S$ along paths that avoid $S\cap T$
is bounded; for a finite graph, iff no closed walk avoiding $S\cap T$ meets
both $S$ and $T$, i.e.\ no strongly connected component of $G$ with
$S\cap T$ removed meets both, a condition decided in linear time; the
bound is then $\min(|S\setminus T|,|T\setminus S|)$ and the join topology
is $J_{S\cap T}$. The theorem interpolates between two known
facts: on a finite category every pair of topologies has a preserved join
\cite{Finitejoin}, which is the acyclic case, and on the poset
$\mathbb N$ the topologies of the even and of the odd numbers do not
\cite{Topjoin}, which is the infinite path with unbounded alternation. A
third section records the finite action monoids: there the ideal
topologies are indexed by sets of idempotent actions, resets, and the
reading of the closure is ``every future through a reset lands in the
invariant''.

\subsection*{What is known.}
Samenesses are from \cite{Same}; the finite-iteration criterion in a
Grothendieck topos from \cite{Stabjoin}, in exact categories with closure
operators, including nuclei on a frame as the poset case, from
\cite{Fijoin}; the finite-category join theorem from \cite{Finitejoin};
the counterexample on $\mathbb N$ from \cite{Topjoin}. Abstract domains as
closure operators, the lattice of abstract domains, the reduced product
and the least common abstraction are in \cite{CC79,Ward,GR97,GRS00}, and
the computation of the join of two closure operators by transfinite
alternation (Lemma~\ref{lem:iter}) is standard and is recalled only to fix
notation; the
iterated reduction computing the reduced product on a product domain is in
\cite{Granger,CCM11}; cut sets of control-flow graphs in
\cite{Bourdoncle}. Lawvere--Tierney topologies, sheafification and the
bidense description of the maps it inverts are in \cite{MLM,Elephant}.

\subsection*{What is new.}
Theorem~\ref{thm:lattice} (closure operators on a complete lattice without
meet preservation: meet, join as alternating chains plus one cancellation,
criterion, necessity of cancellation),
Proposition~\ref{prop:reach} (reachability closures) and
Theorem~\ref{thm:cut} (cut-point topologies on the path category of a
graph, with the closed-walk criterion, its linear-time test and the
identification of the join topology), together with the dictionary of Section~\ref{sec:monoid}. The
engine is the criterion of \cite{Fijoin}; what is added is its application
to closures that are not nuclei and to the object-generated topologies of
infinite path categories, where the combinatorics of alternation along
paths decides the question.

\section{Closure operators on a complete lattice}\label{sec:lattice}

Let $L$ be a complete lattice, regarded as a category with a morphism
$x\to y$ iff $x\le y$. A \emph{closure operator} on $L$ is a map
$\rho\colon L\to L$ that is extensive ($x\le\rho x$), monotone and
idempotent; its set of fixed points $\Fix\rho=\rho(L)$ is closed under
arbitrary meets, and $\rho x$ is the least fixed point above $x$. The
closure operators on $L$ form a complete lattice $\uco(L)$ under the
pointwise order \cite{Ward}, which is the reverse inclusion of fixed-point
sets; the meet of $\rho,\rho'$ is $x\mapsto\rho x\wedge\rho'x$, with fixed
points the meets $a\wedge b$, $a\in\Fix\rho$, $b\in\Fix\rho'$, and the join
$\rho\sqcup\rho'$ is the closure operator with fixed points
$\Fix\rho\cap\Fix\rho'$. In abstract interpretation \cite{CC79,GR97},
$\rho\le\rho'$ means that $\rho$ is more precise, the meet is the reduced
product and the join the least common abstraction.

\begin{lem}[standard]\label{lem:iter}
For closure operators $\rho,\rho'$ and $x\in L$ define $x_0=x$,
$x_{\alpha+1}=\rho\rho'x_\alpha$ and $x_\lambda=\bigvee_{\alpha<\lambda}
x_\alpha$ at limits. The sequence is increasing, stabilises at some
ordinal, and its value is $(\rho\sqcup\rho')x$.
\end{lem}
\begin{proof}
If $x_{\alpha+1}=x_\alpha=:z$ then $\rho'z\le\rho\rho'z=z$ gives $\rho'z=z$
and then $\rho z=\rho\rho'z=z$, so $z\in\Fix\rho\cap\Fix\rho'$; every common
fixed point above $x$ is above all $x_\alpha$, so $z$ is the least one.
\end{proof}

\begin{defi}
For a closure operator $\rho$ on $L$ put
\[
W_\rho:=\{x\le y:\ \rho x=\rho y\}=\{x\le y:\ y\le\rho x\},
\]
the inequalities collapsed by $\rho$, i.e.\ the refinements $y$ of $x$
that $\rho$ does not distinguish from $x$.
\end{defi}

$W_\rho$ is a sameness on the poset $L$: it contains the identities, is
closed under composition ($x\le y\le z$ with $z\le\rho y\le\rho\rho x=\rho
x$), and under both cancellations ($x\le y\le z$ with $z\le\rho x$ gives
$y\le\rho x$ and $z\le\rho x\le\rho y$). No meet preservation is used. If
$\rho\le\rho'$ then $W_\rho\subseteq W_{\rho'}$: a less precise analysis
collapses more.

\begin{thm}\label{thm:lattice}
Let $\rho,\rho'$ be closure operators on a complete lattice $L$.
\begin{enumerate}
\item[(i)] $W_\rho\cap W_{\rho'}=W_{\rho\sqcap\rho'}$: the reduced product
distinguishes $x\le y$ iff one of the two abstractions does.
\item[(ii)] $W_\rho\vee W_{\rho'}=\{x\le y:\ y\le(\rho\rho')^kx\text{ for
some }k\ge0\}$ in $\Same(L)$; for such a pair, with $z=(\rho\rho')^kx$,
both $x\le z$ and $y\le z$ are alternating chains of steps in $W_{\rho'}$
and $W_\rho$, so every element of the join is obtained from two alternating
chains by one cancellation.
\item[(iii)] $W_\rho\vee W_{\rho'}=W_{\rho\sqcup\rho'}$ iff for every $x\in
L$ there is $k$ with $(\rho\rho')^kx=(\rho\sqcup\rho')x$: the alternation
of Lemma~\ref{lem:iter} stops at a finite stage at every point.
\item[(iv)] $\rho\rho'=\rho\sqcup\rho'$ (one alternation suffices
uniformly) iff $\rho\rho'$ is idempotent iff $\rho$ maps $\Fix\rho'$ into
$\Fix\rho'$.
\item[(v)] The compositional closure of $W_\rho\cup W_{\rho'}$, the class of
alternating chains, can be strictly smaller than the join; it equals the
join when $\rho$ and $\rho'$ preserve binary meets.
\end{enumerate}
\end{thm}
\begin{proof}
(i) $x\le y$ lies in $W_{\rho\sqcap\rho'}$ iff $y\le\rho x\wedge\rho'x$ iff
$y\le\rho x$ and $y\le\rho'x$.

(ii) Let $D$ be the right-hand side. $D$ is a sameness: if $y\le(\rho
\rho')^kx$ and $z\le(\rho\rho')^ly$ then $z\le(\rho\rho')^{k+l}x$; if
$x\le y\le z$ and $z\le(\rho\rho')^kx$ then $y\le(\rho\rho')^kx$ and
$z\le(\rho\rho')^kx\le(\rho\rho')^ky$. It contains $W_\rho$ and
$W_{\rho'}$ ($y\le\rho x\le\rho\rho'x$, $y\le\rho'x\le\rho\rho'x$). So the
join $J$ of $W_\rho$ and $W_{\rho'}$ is contained in $D$. Conversely let
$x\le y\le z:=(\rho\rho')^kx$, and put $x_0=x$, $x_{2i+1}=\rho'x_{2i}$,
$x_{2i+2}=\rho x_{2i+1}$, so that $x_{2k}=z$. The chain $x=x_0\le
x_1\le\cdots\le x_{2k}=z$ has its steps alternately in $W_{\rho'}$ and
$W_\rho$, so $(x,z)\in J$. The chain $y=y\vee x_0\le y\vee x_1\le\cdots\le
y\vee x_{2k}=z$ (the last equality as $y\le z$) has the same property: if
$x_{j+1}=dx_j$ with $d\in\{\rho,\rho'\}$, then $y\vee x_{j+1}\le d(y\vee
x_j)$, both $y$ and $dx_j$ being below $d(y\vee x_j)$; so $(y,z)\in J$.
Then $(x,y)\in J$ by cancellation from $(x,z)$ and $(y,z)$.

(iii) $W_{\rho\sqcup\rho'}=\{x\le y:y\le(\rho\sqcup\rho')x\}$ contains $D$,
since $(\rho\rho')^kx\le(\rho\sqcup\rho')x$. If every fibre top
$(\rho\sqcup\rho')x$ is some $(\rho\rho')^kx$ the two classes agree;
if for some $x$ no $(\rho\rho')^kx$ equals $(\rho\sqcup\rho')x$, then
$(x,(\rho\sqcup\rho')x)$ lies in $W_{\rho\sqcup\rho'}$ but not in $D$,
as $(\rho\sqcup\rho')x\le(\rho\rho')^kx$ would force equality.

(iv) If $\rho\rho'$ is idempotent and $a\in\Fix\rho'$, then
$\rho'\rho a\le\rho\rho'\rho a=\rho\rho'\rho\rho'a=\rho\rho'a=\rho a$, so
$\rho a\in\Fix\rho'$. If $\rho$ preserves $\Fix\rho'$ then
$\rho\rho'x\in\Fix\rho'$, so $\rho\rho'\rho\rho'x=\rho\rho\rho'x=\rho\rho'x$.
An idempotent $\rho\rho'$ is a closure operator above $\rho$ and $\rho'$ and
below every closure operator above both, so it is $\rho\sqcup\rho'$;
conversely $\rho\sqcup\rho'$ is idempotent.

(v) Example~\ref{ex:square} gives a strict inclusion. If $\rho,\rho'$
preserve binary meets, let $x\le y\le z=x_{2k}$ with the chain $x_j$ of
(ii) and put $u_j=y\wedge x_j$, so $u_0=x$, $u_{2k}=y$ and $u_j\le
u_{j+1}$. If $x_{j+1}=dx_j$ then $u_{j+1}=y\wedge dx_j\le dy\wedge
dx_j=d(y\wedge x_j)=du_j$, so $u_j\le u_{j+1}$ lies in $W_d$ and
$x=u_0\le\cdots\le u_{2k}=y$ is an alternating chain. (When $L$ is a
frame this is the poset case of Theorem~A of \cite{Fijoin}.)
\end{proof}

\begin{ex}\label{ex:square}
Let $L=\{0<1,2<3\}$ be the four-element Boolean lattice, $\rho$ the closure
with $\Fix\rho=\{0,3\}$ ($\rho1=\rho2=3$) and $\rho'$ the closure with
$\Fix\rho'=\{1,3\}$ ($\rho'0=1$, $\rho'2=3$). Then $W_\rho=\{1\le3,2\le3\}$,
$W_{\rho'}=\{0\le1,2\le3\}$ (plus identities), $\rho\rho'0=3$, so
$0\le2$ lies in the join by (ii); but the only non-identity step out of
$0$ in $W_\rho\cup W_{\rho'}$ is $0\le1$, and $1\not\le2$, so $0\le2$ is
not an alternating chain: it is obtained from $0\le3$ and $2\le3$ by
cancellation. Here $\rho(1\wedge2)=0\ne3=\rho1\wedge\rho2$.
\[
\begin{tikzcd}[row sep=small,column sep=small]
& 3 & \\
1\arrow[ur,"W_\rho"] & & 2\arrow[ul,"W_\rho\cap W_{\rho'}"'] \\
& 0\arrow[ul,"W_{\rho'}"]\arrow[ur,dashed,"\text{join only}"'] &
\end{tikzcd}
\]
\end{ex}

\begin{rem}
In a finite lattice the alternation of Lemma~\ref{lem:iter} stops after
at most $|L|$ steps, so (iii) always holds there and the content of the
theorem is (ii) and (v): the join is generated by alternating chains
\emph{and} cancellation. On an infinite lattice (iii) can fail; the
simplest instance is in Section~\ref{sec:reach}.
\end{rem}

\section{Reachability closures}\label{sec:reach}

Let $X$ be a set of states, $R\subseteq X\times X$ a transition relation
and $I\subseteq X$ a set of initial states. On the complete lattice
$L=\cP(X)$ put
\[
\rho_{R,I}(S):=R^*(S\cup I),
\]
the set of states reachable from $S$ or from $I$ along $R$: the collecting
semantics of ``run $R$ from $S$''.

\begin{prop}\label{prop:reach}
\begin{enumerate}
\item[(i)] $\rho_{R,I}$ is a closure operator on $\cP(X)$ with fixed points
the $R$-closed sets containing $I$ (the inductive invariants); it
preserves binary intersections only in degenerate cases.
\item[(ii)] $\rho_{R,I}\sqcap\rho_{R',I'}$ is $S\mapsto R^*(S\cup I)\cap
R'^*(S\cup I')$ and $\rho_{R,I}\sqcup\rho_{R',I'}=\rho_{R\cup R',I\cup
I'}$; for $k\ge1$, $(\rho_{R,I}\rho_{R',I'})^kS$ is the set of states
reachable from $S\cup I'$ by a path of at most $k$ blocks, each an
$R'$-path followed by an $R$-path, together with the states reachable from
$I$ by an $R$-path followed by at most $k-1$ such blocks (the initial
states of the first system enter only after the $R'$-phase of a block).
\item[(iii)] $W_{\rho_{R,I}}\vee W_{\rho_{R',I'}}$ is the class of pairs
$S\subseteq T$ with $T\subseteq(\rho_{R,I}\rho_{R',I'})^kS$ for some $k$,
i.e.\ every state of $T$ is reachable in the sense of (ii) with a bounded
number of blocks, and it equals
$W_{\rho_{R\cup R',I\cup I'}}$ iff the alternation depth of $(R\cup
R')$-reachability is finite from every $S$. If $X$ is finite this holds
with at most $|X|$ blocks.
\item[(iv)] On $X=\mathbb N$ with $R=\{(2i,2i+1)\}$, $R'=\{(2i+1,2i+2)\}$
and $I=I'=\emptyset$, $(\rho_R\rho_{R'})^k\{0\}=[0,2k-1]$ for $k\ge1$
(the first $R'$-step from $0$ is empty), which never reaches
$\rho_{R\cup R'}\{0\}=\mathbb N$, so $W_{\rho_R}\vee W_{\rho_{R'}}\subsetneq
W_{\rho_{R\cup R'}}$.
\end{enumerate}
\end{prop}
\begin{proof}
(i) Extensive, monotone and idempotent since $R^*$ is reflexive and
transitive; the fixed points are the sets closed under $R$ and containing
$I$. For $x\xrightarrow{R}z\xleftarrow{R}y$ with $x\ne y$ one has
$z\in\rho\{x\}\cap\rho\{y\}$ but $\rho(\{x\}\cap\{y\})=\rho\emptyset=R^*I$.
(ii) The meet is pointwise; the common fixed points of the two closures
are the sets closed under $R$ and $R'$ containing $I\cup I'$, the fixed
points of $\rho_{R\cup R',I\cup I'}$, which is therefore the join. For the
iterates, $\rho_{R,I}\rho_{R',I'}(A)=R^*R'^*(A\cup I')\cup R^*I$; writing
$B_k(A)$ for the states reachable from $A$ by at most $k$ blocks, so that
$B_1B_k=B_{k+1}$ and $B_k(A)\subseteq B_{k+1}(A)$ as blocks may be empty,
induction on $k$ gives $(\rho_{R,I}\rho_{R',I'})^kS=B_k(S\cup I')\cup
B_{k-1}(R^*I)$. (iii) is Theorem~\ref{thm:lattice}(ii),
(iii); if $X$ is finite, every block that changes the set adds a state,
so the iteration of Lemma~\ref{lem:iter} stops after at most $|X|$ blocks.
(iv) is a direct computation, and Theorem~\ref{thm:lattice}(iii).
\end{proof}

\begin{rem}
Example~\ref{ex:square} is of this form: $X=\{p,q\}$, $R=\{(p,q),(q,p)\}$,
$I=\emptyset$ and $R'=\emptyset$, $I'=\{p\}$. The pair
$\emptyset\subseteq\{q\}$ is collapsed by the join (from $\emptyset$ the
second system reaches $p$ and the first then $q$) but by no alternating
chain of collapses, the only refinement of $\emptyset$ collapsed by one of
the two systems being $\{p\}\not\subseteq\{q\}$.
\end{rem}

\begin{rem}\label{rem:granger}
The iterated reduction of Granger \cite{Granger} and of Cousot, Cousot and
Mauborgne \cite{CCM11}, which alternately refines the two components of a
pair $(a,b)$ in a product of abstract domains and is the Nelson--Oppen
procedure when the domains are logical theories, computes the reduced
product, the \emph{meet} $\rho\sqcap\rho'$, by a decreasing iteration on
the product domain. The alternation $\rho\rho'\rho\rho'\cdots$ of
Lemma~\ref{lem:iter} computes the \emph{join}, the least common
abstraction, by an increasing iteration on $L$. The two are not to be
confused: the first terminates in one step in the sense of (i) of
Theorem~\ref{thm:lattice}, the second is governed by (iii).
\end{rem}

\section{Cut-point topologies on the path category of a graph}\label{sec:cut}

Let $G$ be a directed graph, finite or infinite, with vertex set $V$, and
let $\cP(G)$ be its path category: objects the vertices, morphisms the
finite paths, identities the empty paths. A presheaf $X$ on $\cP(G)$ is a
set $X(c)$ for each vertex and a restriction $X(c)\to X(d)$, $x\mapsto
x\cdot f$, for each path $f\colon d\to c$, a flow of sets along the control
flow; a subpresheaf is a family of subsets closed under restriction. Write
$\cE=\Psh(\cP(G))$, a Grothendieck topos.

For a set $S\subseteq V$ of vertices, cut points, let $D_S$ be the set of
paths that pass through a vertex of $S$, the empty paths at vertices of $S$
included. $D_S$ is a two-sided ideal of $\cP(G)$ (closed under composition
with arbitrary paths on either side) and idempotent, $D_S\circ D_S=D_S$,
because $f=f\circ\mathrm{id}_s$ for a path $f$ through $s\in S$. For a
vertex $c$ let $D_S(c)$ be the sieve of paths into $c$ through $S$.

\begin{lem}\label{lem:JS}
The sieves containing $D_S(c)$, for all $c$, form a Grothendieck topology
$J_S$ on $\cP(G)$. Its sheaves are the presheaves $X$ with $X\cong
\Ran_i(X|_S)$ for the inclusion $i$ of the full subcategory of $\cP(G)$ on
the vertices of $S$, so that restriction gives $\Sh(J_S)\simeq\Psh(S)$,
where $\Psh(S)$ denotes the presheaves on that full subcategory (its
morphisms are all paths of $G$ between vertices of $S$, not only the edges
inside $S$), and its closure operator on
subpresheaves $A\subseteq X$ is
\[
c_S(A)(c)=\{x\in X(c):\ x\cdot f\in A\text{ for every path }f\colon s\to c
\text{ with }s\in S\}.
\]
For two sets $S,T$ one has $c_Sc_T=c_{D_S\circ D_T}$, where
$c_D(A)(c)=\{x:x\cdot f\in A\ \forall f\in D(c)\}$ for an ideal $D$; so
$(c_Sc_T)^k=c_{(D_S\circ D_T)^k}$, and $D_S\circ D_T$ consists of the paths
with a visit to $T$ followed by a visit to $S$.
\end{lem}
\begin{proof}
Stability: the pullback of $D_S(c)$ along $g\colon e\to c$ is $\{h:gh\in
D_S(c)\}\supseteq D_S(e)$. Transitivity: let $R$ be a sieve on $c$ with
$h^*R\supseteq D_S(\mathrm{dom}\,h)$ for all $h\in D_S(c)$; for
$f=u\circ v\in D_S(c)$ with $u\colon s\to c$, $s\in S$, we have $u\in
D_S(c)$ and $\mathrm{id}_s\in D_S(s)\subseteq u^*R$, so $u\in R$ and
$f\in R$. Restriction $\Psh(\cP(G))\to\Psh(S)$ is left exact with right
adjoint $\Ran_i$, and $\Ran_i$ is fully faithful since $i$ is, so
restriction is the inverse image of a geometric embedding; a sieve
$R\hookrightarrow y(c)$ is inverted by restriction iff $R$ contains every
path from a vertex of $S$ into $c$, i.e.\ iff $R\supseteq D_S(c)$, so the
embedding is $\Sh(J_S)$ \cite[A4.3, A4.5]{Elephant}. The closure of $A$
for $J_S$ consists of the $x$ whose sieve $\{f:x\cdot f\in A\}$ is
covering, which is the displayed formula; and $c_Dc_E(A)=\{x:x\cdot f\in
c_E(A)\ \forall f\in D\}=\{x:x\cdot f\cdot g\in A\ \forall f\in D,g\in E\}
=c_{D\circ E}(A)$, with $f\circ g$ the path $g$ followed by $f$.
\end{proof}

\begin{lem}\label{lem:alt}
A path $f$ lies in $(D_S\circ D_T)^k$ iff there are vertices
$t_1,s_1,\dots,t_k,s_k$ occurring along $f$ in this order, positions
non-decreasing, with $t_i\in T$ and $s_i\in S$. A path through a vertex of
$S\cap T$ lies in $(D_S\circ D_T)^k$ for every $k$. Hence the chain
$(D_S\circ D_T)^k$ stabilises iff the number of alternations ``a visit to
$T$, then a visit to $S$'' along paths avoiding $S\cap T$ is bounded, and
then $(D_S\circ D_T)^{a+1}=D_{S\cap T}$ for the bound $a$.
\end{lem}
\begin{proof}
$f\in(D_S\circ D_T)^k$ iff $f=u_kv_k\cdots u_1v_1$ with $v_i\in D_T$, $u_i\in
D_S$, the splitting points being arbitrary vertices along $f$ and the
factors possibly empty; this is the stated condition. At a common vertex
$x$ all factors may be taken empty. If the alternation along paths
avoiding $S\cap T$ is bounded by $a$, a path in level $a+1$ passes through
$S\cap T$, and conversely such paths are in every level; if it is
unbounded, the chain drops at infinitely many levels.
\end{proof}

\begin{thm}\label{thm:cut}
Let $G$ be a directed graph and $S,T\subseteq V$. In
$\Same(\Psh(\cP(G)))$,
\[
W_{J_S}\vee W_{J_T}=W_{J_S\vee J_T}
\]
iff the number of alternations between visits to $T$ and visits to $S$
along paths of $G$ avoiding $S\cap T$ is bounded, and then $J_S\vee
J_T=J_{S\cap T}$, $\Sh(J_S)\cap\Sh(J_T)\simeq\Psh(S\cap T)$. If $G$ is
finite this holds iff no closed walk of $G$ avoiding $S\cap T$ meets both
$S$ and $T$, i.e.\ iff no strongly connected component of the induced
subgraph $G-(S\cap T)$ contains both a vertex of $S$ and a vertex of $T$;
the bound is then at most $\min(|S\setminus T|,|T\setminus S|)$, and the
condition is decided in time $O(|V|+|E|)$ by a decomposition into strongly
connected components. In particular, for disjoint $S,T$ on a finite graph
the join is preserved iff no strongly connected component of $G$ meets both
$S$ and $T$.
\end{thm}
\begin{proof}
Let $c=c_S$, $c'=c_T$ and let $\bar c$ be the least closure operator
above both, the closure operator of the join topology; by
\cite[Thm.~B]{Fijoin}, $W_{J_S}\vee W_{J_T}=W_{\bar c}$ iff $(cc')^k=\bar c$
for some $k$. If the alternation is bounded by $a$, then by
Lemmas~\ref{lem:JS} and~\ref{lem:alt} $(cc')^{a+1}=c_{D_{S\cap T}}$ is
idempotent, hence a closure operator above $c$ and $c'$, and every closure
operator above both is above $cc'$ and so above $(cc')^{a+1}$; thus $\bar
c=(cc')^{a+1}=c_{S\cap T}$, the criterion holds, and the join topology is
$J_{S\cap T}$. If the alternation is unbounded, the chain of ideals
$(D_S\circ D_T)^k$ does not stabilise by Lemma~\ref{lem:alt}, and an ideal
$D$ is determined by $c_D$: for $X=y(c)$, $D(c)$ is the least sieve $A$ on
$c$ with $\mathrm{id}_c\in c_D(A)(c)$. Since $c,c'\le\bar c$ and $\bar c$ is
idempotent, $(cc')^k\le\bar c$ for all $k$ and $c\bar c=c'\bar c=\bar c$.
If $\bar c=(cc')^k$ for some $k$, then $(cc')^m=(cc')^{m-k}\bar c=\bar
c=(cc')^k$ for all $m\ge k$, so $(D_S\circ D_T)^m=(D_S\circ D_T)^k$ for all
$m\ge k$ and the chain stabilises, a contradiction; so the criterion fails.
For finite $G$: if a closed walk avoiding $S\cap T$ meets $s\in S$ and
$t\in T$, going round it $m$ times gives alternation $\ge m$; conversely a
path avoiding $S\cap T$ with alternation $\ge|S\setminus T|+1$ repeats an
$S$-vertex, $s_i=s_j$ with $i<j$, with a $T$-visit $t_{i+1}$ strictly
between (a visit to $S$ and one to $T$ at the same position would be a
vertex of $S\cap T$), and the subwalk from $s_i$ to $s_j$ is a closed walk
avoiding $S\cap T$ meeting $S$ and $T$; likewise with $T\setminus S$ and
the $t_i$. Two vertices lie on a common closed walk of $G-(S\cap T)$ iff
they lie in the same strongly connected component, and the components are
computed in linear time \cite{Tarjan}.
\end{proof}

\begin{ex}\label{ex:graphs}
(a) $G$: $s\to t\to s'$, $S=\{s,s'\}$, $T=\{t\}$. The alternation is at most
one, the join is preserved, and $J_S\vee J_T=J_\emptyset$ is the degenerate
topology: a $J_S$-sheaf has $X(t)\cong X(s)$, a $J_T$-sheaf has
$X(s)=1$ and $X(s')\cong X(t)$, so a common sheaf is terminal.
(b) $G$: the two-cycle $s\rightleftarrows t$, $S=\{s\}$, $T=\{t\}$. Going
round the cycle gives unbounded alternation, and $W_{J_S}\vee
W_{J_T}\subsetneq W_{J_S\vee J_T}$; adding the vertex $s$ to $T$ restores
preservation, as every closed walk then meets $S\cap T=\{s\}$.
(b$'$) $G$: $s\rightleftarrows u\rightleftarrows t$, $S=\{s\}$, $T=\{t\}$.
The simple cycles $s\to u\to s$ and $u\to t\to u$ each meet only one of
the two sets, but the closed walk $s\to u\to t\to u\to s$ meets both,
the alternation is unbounded and the join is not preserved: the condition
of Theorem~\ref{thm:cut} is on closed walks, equivalently on strongly
connected components, not on simple cycles.
(c) $G$ the infinite path $0\to1\to2\to\cdots$, $S$ the even and $T$ the
odd vertices: no cycles, unbounded alternation, and the failure of
\cite{Topjoin}.
\[
\begin{tikzcd}[column sep=small]
s\arrow[r] & t\arrow[r] & s'
\end{tikzcd}
\qquad
\begin{tikzcd}[column sep=small]
s\arrow[r,bend left] & t\arrow[l,bend left]
\end{tikzcd}
\qquad
\begin{tikzcd}[column sep=small]
s\arrow[r,bend left] & u\arrow[l,bend left]\arrow[r,bend left] &
t\arrow[l,bend left]
\end{tikzcd}
\qquad
\begin{tikzcd}[column sep=small]
0\arrow[r] & 1\arrow[r] & 2\arrow[r] & \cdots
\end{tikzcd}
\]
\end{ex}

\begin{rem}
A finite acyclic $G$ has a finite path category, and \cite{Finitejoin}
gives preservation for all pairs of topologies on a finite category;
Theorem~\ref{thm:cut} locates the loops that break it. In the terms of
\cite{Bourdoncle}, a set of cut points is a set meeting every cycle; the
theorem concerns two such sets and asks that their common part meet every
closed walk carrying points of both; for one set, meeting every simple
cycle and meeting every closed walk are the same condition, for the pair
they are not (Example~\ref{ex:graphs}(b$'$)). The bound
$\min(|S\setminus T|,|T\setminus S|)$ on the alternation is the number of
mutual reconstructions, from the values at $S$ and from the
values at $T$, after which the reconstruction from $S\cap T$ is reached.
\end{rem}

\begin{rem}
A presheaf on $\cP(G)$ restricts against the direction of the edges,
$X(c)\to X(d)$ for a path $d\to c$, so the flows of
Section~\ref{sec:cut} are the backward ones; forward flows are presheaves
on $\cP(G^{\mathrm{op}})$, the graph with the edges reversed. Closed walks
and strongly connected components are the same for $G$ and
$G^{\mathrm{op}}$, so the criterion of Theorem~\ref{thm:cut} and the
bound are the same in both directions.
\end{rem}

\begin{que}
Is $\Sh(J_S)\cap\Sh(J_T)=\Sh(J_{S\cap T})$ for all $S,T$, i.e.\ $J_S\vee
J_T=J_{S\cap T}$ also when the alternation is unbounded? The intersection
of the ideals $(D_S\circ D_T)^k$ is $D_{S\cap T}$ in all cases, and for the
infinite path with $S$ the even and $T$ the odd vertices the join is
degenerate \cite{Topjoin}, which is $J_\emptyset$.
\end{que}

\section{Finite action monoids}\label{sec:monoid}

A deterministic transition system with action alphabet $\Sigma$ is a
coalgebra $X\to X^\Sigma$, i.e.\ a right $\Sigma^*$-set, and the coalgebra
homomorphisms are the equivariant maps; the presheaf topos of right
$M$-sets for a monoid $M$ is the general form. Its \emph{ideal
topologies} are those with a least covering sieve, and they correspond to
the idempotent two-sided ideals $D\subseteq M$, with closure $c_D(A)=\{x:
x\cdot D\subseteq A\}$ \cite{Stabjoin,Fijoin}: a state is in the closure
of the invariant $A$ iff all its $D$-futures lie in $A$.

\begin{lem}\label{lem:idem}
In a finite monoid $M$ a two-sided ideal $D$ is idempotent iff it is
generated by idempotents, $D=M\,E\,M$ with $E=\{e\in D:e^2=e\}$.
\end{lem}
\begin{proof}
$(MeM)(MeM)\ni ue\,ev=uev$ gives $MEM\subseteq(MEM)^2\subseteq MEM$.
Conversely let $D=D^2$, so $D=D^n$ for all $n$, and write $d\in D$ as
$d=s_1\cdots s_n$ with $s_i\in D$ and $n>|M|$. Two prefixes coincide,
$s_1\cdots s_i=s_1\cdots s_j$ with $i<j$, so with $w=s_{i+1}\cdots s_j\in D$
the prefix $p=s_1\cdots s_i$ satisfies $p=pw^m$ for all $m$; some power
$e=w^m$ is idempotent, $e\in D$, and $d=pe\,s_{j+1}\cdots s_n\in MeM$.
\end{proof}

\begin{prop}
On the right $M$-sets of a finite monoid $M$ the ideal topologies are
indexed by the ideals generated by sets $E$ of idempotents, with
$c_{MEM}(A)=\{x:x\,u\,e\in A\text{ for all }u\in M,e\in E\}$: a state is in
the closure of the invariant $A$ iff every future that passes through a
reset $e\in E$ lands in $A$. By \cite{Stabjoin}, $W_D\vee W_{D'}=W_{J_D\vee
J_{D'}}$ for any two of them, with the stabilisation rank of
$(DD')^k$ as the bound on the alternation. For the free monoid $\Sigma^*$,
and more generally for a monoid with a length function additive and
positive on non-identity elements, the only idempotent ideals are
$\emptyset$ and $M$, so the deterministic systems over a free alphabet
have no non-trivial ideal topology.
\end{prop}
\begin{proof}
The first statement is Lemma~\ref{lem:idem} and the description of
$c_D$, as $x\cdot MeM\subseteq A$ iff $x\,u\,e\in A$ for all $u$ (the
right factor preserves $A$). The second is \cite{Stabjoin} together with
finiteness. For the last, a shortest non-identity element of an idempotent
ideal $D\ne M$ is a product of two elements of $D$ of larger length.
\end{proof}

\begin{rem}
Nontrivial ideal topologies on deterministic systems thus need relations
among the actions, typically idempotent actions (resets, normalisation
steps). In a band ($w^2=w$ for all $w$) every ideal is idempotent, and
for the free band on $\Sigma$ the ideals $D_a$ of words containing the
letter $a$ satisfy $D_aD_b=D_a\cap D_b$: one alternation. The contracted
cycle monoids of \cite{Finitejoin}, with stabilisation rank growing with
the cycle, are the instances of unbounded alternation in this setting.
\end{rem}

\section{Questions}

\begin{que}
Which pairs of standard numerical abstract domains (intervals, congruences,
octagons, polyhedra, on $\cP(\mathbb Z^n)$) have unbounded alternation at
some element, so that the join of the two classes of collapsed pairs is
not the class of the least common abstraction? The pairs checked
(intervals and congruences) stabilise in two steps because the interval
hull of an infinite congruence class is a product of lines and points.
\end{que}

\begin{que}
Extend Theorem~\ref{thm:cut} to flows with transfer functions, i.e.\ to
presheaves on a quotient of the path category by relations among the
edges (idempotent edges, commuting edges), where the ideals $D_S$ are
replaced by their images; for which relations does the cycle criterion
survive?
\end{que}

\section*{Funding}
No funding was received for this work.

\section*{Declarations}
The author has no competing interests to declare. The finite checks
mentioned in the text (closure operators on small lattices, random
transition systems, the cycle criterion on random graphs) are the files
\texttt{uco\_same.py}, \texttt{reach\_same.py} and \texttt{cutpoints.py}
in the directory \texttt{scripts/comparison} of the repository
\url{https://github.com/r-shibusawa/cubical-strict-j}, archived on Zenodo
under the concept DOI 10.5281/zenodo.21405961.

\begin{thebibliography}{9}
\bibitem{Bourdoncle} F.~Bourdoncle, \emph{Efficient chaotic iteration
strategies with widenings}, in: Formal Methods in Programming and their
Applications, Lecture Notes in Comput. Sci. 735, Springer, 1993, 128--141.
\bibitem{CC79} P.~Cousot, R.~Cousot, \emph{Systematic design of program
analysis frameworks}, in: Proc. 6th ACM Symposium on Principles of
Programming Languages (POPL 1979), 269--282.
\bibitem{CCM11} P.~Cousot, R.~Cousot, L.~Mauborgne, \emph{The reduced
product of abstract domains and the combination of decision procedures},
in: FoSSaCS 2011, Lecture Notes in Comput. Sci. 6604, Springer, 2011,
456--472.
\bibitem{DT} D.~Dikranjan, W.~Tholen, \emph{Categorical Structure of Closure
Operators}, Mathematics and its Applications 346, Kluwer, 1995.
\bibitem{GR97} R.~Giacobazzi, F.~Ranzato, \emph{Refining and compressing
abstract domains}, in: ICALP 1997, Lecture Notes in Comput. Sci. 1256,
Springer, 1997, 771--781.
\bibitem{GRS00} R.~Giacobazzi, F.~Ranzato, F.~Scozzari, \emph{Making
abstract interpretations complete}, J. ACM \textbf{47} (2000), 361--416.
\bibitem{Granger} P.~Granger, \emph{Improving the results of static analyses
of programs by local decreasing iterations}, in: FSTTCS 1992, Lecture Notes
in Comput. Sci. 652, Springer, 1992, 68--79.
\bibitem{Elephant} P.~T.~Johnstone, \emph{Sketches of an Elephant: A Topos
Theory Compendium}, Oxford Logic Guides 43--44, Oxford University Press, 2002.
\bibitem{MLM} S.~Mac Lane, I.~Moerdijk, \emph{Sheaves in Geometry and Logic},
Springer, 1992.
\bibitem{Same} R.~Shibusawa, \emph{The sameness lattice of a category:
weak-equivalence structures, Galois functoriality, and comparison intervals},
preprint (2026), archived at DOI 10.5281/zenodo.22823772.
\bibitem{Topjoin} R.~Shibusawa, \emph{Meets and joins of Grothendieck
topologies in the sameness lattice}, preprint (2026), archived at DOI
10.5281/zenodo.22978446.
\bibitem{Finitejoin} R.~Shibusawa, \emph{Grothendieck topologies on a finite
category form a sublattice of the lattice of two-out-of-three classes},
preprint (2026), archived at DOI 10.5281/zenodo.23204943.
\bibitem{Stabjoin} R.~Shibusawa, \emph{Joins of Lawvere--Tierney topologies
in the lattice of two-out-of-three classes: the finite-iteration criterion},
preprint (2026), archived at DOI 10.5281/zenodo.23487202.
\bibitem{Fijoin} R.~Shibusawa, \emph{Joins of localisations in the lattice of
two-out-of-three classes: the finite-iteration criterion for closure
operators on exact categories, with applications to toposes, torsion
theories and fuzzy sets}, preprint (2026), archived at DOI
10.5281/zenodo.23691896.
\bibitem{Tarjan} R.~E.~Tarjan, \emph{Depth-first search and linear graph
algorithms}, SIAM J. Comput. \textbf{1} (1972), 146--160.
\bibitem{Ward} M.~Ward, \emph{The closure operators of a lattice}, Ann. of
Math. (2) \textbf{43} (1942), 191--196.
\end{thebibliography}

\end{document}
